% \documentclass[fleqn,10pt]{wlscirep}
% fleqn : left alignment of equation

\documentclass[fleqn, 10pt]{wlscirep}
\usepackage[utf8]{inputenc}
\usepackage{amsmath}
\usepackage{amsbsy}
\usepackage{xcolor}
\usepackage[T1]{fontenc}
\usepackage{soul} % 여러 라인에 걸친 밑줄을 사용하기 위한 패키지

% \usepackage[style=authoryear,backend=biber]{biblatex}
\DeclareMathOperator*{\argmin}{arg\,min} % Jan Hlavacek

%from here, code for line numbers
\usepackage{lineno}
\linenumbers

\makeatletter
\patchcmd\@maketitle{%
        {%
        \noindent
        \colorbox{color2}{%
            \parbox{\dimexpr\linewidth-2\fboxsep\relax}{%
            \sffamily\small\textbf\\\theabstract
            }%
        }%
        }%
    }{%
        \nolinenumbers{%
        \noindent
        \colorbox{color2}{%
            \parbox{\dimexpr\linewidth-2\fboxsep\relax}{%
            \internallinenumbers \sffamily\small\theabstract
            }%
        }%
        }%
    }
    {\wlog{true}}{\wlog{false}}
    
\makeatother
%end of code for line numbers

\title{Backward Cumulative Channel Sum-Based Spectral Processing for Radionuclide Identification in Plastic Scintillation Detectors}
% \linenumbers
\author[1]{Ji-Yong So}
\author[2]{Byoungil Jeon}
\author[1,*]{Myungkook Moon}
\affil[1]{Korea Atomic Energy Research Institute, Neutron Science Division, Daejeon, 34057, Korea}
\affil[2]{Korea Atomic Energy Research Institute, Applied Artificial Intelligence Section, Daejeon, 34057, Korea}

\affil[*]{moonmk@kaeri.re.kr}

% \affil[+]{These authors equally contributed to this work}

\keywords{radiation portal monitor, radionuclide identification}

\begin{abstract}

Plastic scintillation detectors are widely used in radiation portal monitors because of their low cost, robustness, and high detection efficiency, but their poor energy resolution limits radionuclide identification under weak-signal and high-background conditions. We propose a simple spectral-processing framework based on the Backward Cumulative Channel Sum (BCCS) transformation for improving radionuclide separability in low-resolution plastic scintillation spectra. The method combines BCCS with two background-compensation schemes, direct subtraction and normalization to background, and applies non-negative least-squares decomposition to estimate radionuclide contributions. Reference measurements of Am-241, Cs-137, and Co-60, together with nine mixed-source configurations acquired under representative radiation portal monitoring conditions, were used for validation. The BCCS transformation produced smooth cumulative profiles that suppressed high-channel fluctuations while preserving source-dependent spectral structure. Both compensation methods enhanced radionuclide-specific signatures and supported stable decomposition of mixed-source spectra. The inferred contributions decreased with increasing source-to-detector distance and approached background-derived thresholds under very weak-signal conditions. These results demonstrate that BCCS-based preprocessing provides an interpretable and computationally lightweight approach for radionuclide identification using existing plastic scintillation detector systems without hardware modification or data-intensive training.

\end{abstract} 

\begin{document}

\flushbottom
\maketitle

\thispagestyle{empty}


\section*{Introduction}

Radiation portal monitors (RPMs) are widely deployed at border crossings, transportation facilities, and industrial sites to detect the illicit movement of radioactive materials.\cite{Sicilliano2005Comparision, Bertrand2014Current,Ely2006TheUse}. Most RPM systems employ large-area plastic scintillation detectors because they offer low cost, mechanical robustness, fast response, and high detection efficiency\cite{Hevener2013EWindow, Lee2020Radioisotope, Paff2017Spectral}. However, plastic scintillators inherently provide poor energy resolution, and the resulting spectra are typically dominated by broad Compton continua rather than well-defined photopeaks. This makes radionuclide identification difficult under realistic field conditions, where net count rates may be low and background radiation may vary over time 

A range of signal-processing approaches has been investigated to improve spectral discrimination in low-resolution detectors. Conventional methods include smoothing and denoising filters6, as well as feature-extraction approaches such as principal component analysis\cite{Runkle2006Analysis,Kim2019InvMat}.More recently, machine-learning and deep-learning methods have been introduced for isotope identification using low-resolution gamma-ray spectra\cite{Kim2019ANN, Lee2023Radionuclide,Jeon2021CompML}. Although these methods can improve classification performance, they often require extensive training datasets, careful tuning, and validation across different operating environments. In addition,computationally intensive inference may not be ideal for embedded systems operating under real-time constraints.

In practical RPM operation, radioactive sources are frequently measured at relatively large source-to-detector distances or under partially shielded conditions, leading to weak net signals. Under such circumstances, spectral differences between radionuclides can be masked by background fluctuations. A robust preprocessing method that enhances spectral separability without requiring high-resolution detectors, elaborate calibration procedures, or large training datasets is therefore highly desirable.

In this study, we propose the Backward Cumulative Channel Sum (BCCS) as an effective spectral transformation for improving radionuclide separability in plastic scintillation detectors. By accumulating counts from high to low channels, the BCCS transformation suppresses fluctuations in the high-channel region while preserving the global spectral tendency of each source. We further investigate two complementary background-compensation strategies—direct subtraction and normalization to background—to isolate source-dependent signatures under low-count conditions. Because the processed spectra retain linear compositional structure, they can be analyzed using non-negative least-squares (NNLS) spectral decomposition. Using reference spectra of Am-241, Cs-137, and Co-60, together with multiple mixed-source measurements acquired under representative RPM conditions, we demonstrate that the proposed framework supports both qualitative radionuclide identification and quantitative estimation of radionuclide contributions using only low-resolution plastic scintillation detectors.

\section*{Results}

\subsection*{Transformation of reference radionuclide spectra}

Figure \ref{fig:bkgref}(a) shows the raw spectra of background radiation and three reference radionuclides---Am-241, Cs-137, and Co-60---measured using a large-volume plastic scintillation detector. As expected for plastic scintillators, all spectra are dominated by broad Compton continua, with most counts concentrated in the low-channel region. Cs-137 and Co-60 show noticeable deviations from background over a relatively broad channel range, whereas Am-241 differs from background primarily in the lowest channels, making reliable identification more difficult under low-count conditions.

Applying the BCCS transformation produces smooth, monotonically decreasing profiles, as shown in Figure \ref{fig:bkgref}(b). This transformation suppresses high-channel statistical fluctuations while preserving the overall spectral structure of each radionuclide. Figures \ref{fig:bkgref}(c) and \ref{fig:bkgref}(d) present the background-subtracted and background-normalized BCCS spectra, respectively. The subtraction-based representation retains the absolute count scale and directly reflects the source contribution relative to background, whereas the normalization-based representation emphasizes relative deviations and improves visual contrast. Both representations provide clearer radionuclide-dependent signatures than the raw spectra.

\subsection*{Background compensation for mixed-source spectra}

Mixed-source measurements were processed using the same two background-compensation strategies. Figure \ref{fig:mixedbccs}(a) shows the background-subtracted BCCS spectra for nine mixed-source configurations composed of Am-241, Cs-137, and Co-60 at various source-to-detector distances. Figure \ref{fig:mixedbccs}(b) shows the corresponding normalization-based spectra.

The two approaches provide complementary views of the same data. Direct subtraction preserves absolute count differences among mixtures, allowing direct comparison of net signal magnitude. In contrast, normalization to background enhances relative spectral contrast and facilitates visual discrimination among mixtures, particularly when the constituent radionuclides contribute in different channel regions or when the overall signal level is sufficiently above background. These observations indicate that both compensation methods are useful, depending on whether the primary goal is absolute quantification or qualitative pattern recognition.

For Am-241, whose spectral contribution is concentrated strongly in the lowest channels, statistical fluctuations may produce small negative values after subtraction. In our analysis, these values are much smaller than the dominant signal amplitude and are attributable to statistical noise rather than physical signal inversion. They were therefore set to zero during the decomposition procedure to prevent unstable fitting.


\subsection*{Spectral decomposition by NNLS fitting}

The processed spectra were further analyzed using NNLS spectral decomposition. Figures \ref{fig:fi_fitbgsub} presents the fitting results for the background-subtracted BCCS spectra, and Figure \ref{fig:fi_fitbgrto} shows the corresponding results for the background-normalized BCCS
spectra. In each panel, the measured spectrum and the fitted reconstruction are shown together with the inferred component contributions from the reference radionuclides.

The fitted spectra reproduce the measured data well in both representations, demonstrating that the processed spectra retain sufficient structure for reliable decomposition. Because meaningful information is concentrated primarily in the low-channel region for plastic scintillation detectors operating under weak-signal conditions, the fitting was restricted to the low-channel portion of the spectrum. This restriction reduces the influence of noise-dominated high-channel data and improves the stability of the inferred coefficients.

The normalization-based spectra often provide a more intuitive visual representation of source presence, especially when one radionuclide produces a stronger relative deviation from background in a limited channel region. Nevertheless, the subtraction-based spectra remain advantageous for preserving physically meaningful count scales. Together, these results
show that both processed representations are compatible with NNLS-based decomposition and can support robust radionuclide identification in mixed-source measurements.

\subsection*{Distance dependence of estimated contributions}

Figure \ref{fig:fi_fitcontrib} summarizes the estimated radionuclide contributions as a function of source-to-detector distance for the subtraction-based and normalization-based analyses. In both cases, the inferred contributions decrease monotonically with increasing distance, consistent with the expected reduction in signal due to geometric attenuation. 

The horizontal dashed lines indicate mean contribution levels estimated from measurements in which the corresponding radionuclide was absent. These values serve as practical background-derived thresholds for interpretation. For Co-60 and Cs-137, the estimated contributions converge toward these threshold bands at large distances, indicating that the corresponding source signals become difficult to distinguish from background under very low-count conditions. Am-241 shows a similar trend, although its behavior is more strongly influenced by the concentration of counts in the lowest channels.

Overall, the results demonstrate that the proposed BCCS-based preprocessing and decomposition framework can identify radionuclide presence and estimate radionuclide contributions in a physically consistent manner, even in weak-signal and mixed-source scenarios.

\section*{Discussion}

The results of this study show that the proposed BCCS-based framework provides a practical approach for enhancing radionuclide identification using low-resolution plastic scintillation detectors. Although such detectors do not provide the spectral resolution needed to resolve photopeaks clearly, the BCCS transformation suppresses high-channel fluctuations while preserving the global source-dependent structure of the spectrum. This enhances spectral separability precisely in the operating regime where plastic scintillators are most often challenged: weak signals, mixed sources, and variable backgrounds.

An important strength of the proposed approach is its simplicity. Unlike many machine-learning methods, the BCCS transformation and the associated background-compensation procedures do not require training datasets, feature engineering, or model retraining when measurement conditions change. Instead, they operate directly on measured spectra, making them suitable for embedded systems and real-time monitoring applications. The cumulative nature of the transformation also reduces the influence of bin-to-bin Poisson variability, which is particularly beneficial when spectra contain only limited net counts. 

The two compensation methods considered here play complementary roles. Direct subtraction preserves absolute source-related count information and is therefore well suited to quantitative decomposition. Normalization to background, on the other hand, emphasizes relative deviations and often improves qualitative visual interpretation, especially in mixed-source cases where subtle spectral differences might otherwise be obscured. The fact that both representations remain compatible with NNLS decomposition is an important practical advantage, as it provides flexibility in implementation and interpretation.

The observed reduction in inferred contributions with increasing source-to-detector distance is consistent with physical expectations and supports the validity of the proposed framework. When the estimated contribution approaches the background-derived threshold, the method effectively indicates the practical detection limit imposed by geometric attenuation and counting statistics rather than by the mathematical method itself. In this sense, the framework provides not only source identification but also an interpretable boundary between detectable and background-indistinguishable conditions.

The proposed method is particularly relevant to field-deployed RPM systems, where low-cost and mechanically robust detectors are often preferred over higher-resolution but more fragile or expensive alternatives such as HPGe systems. Because the method does not require hardware modification and introduces only limited computational overhead, it can be incorporated into existing systems with minimal implementation burden. The general concept may also be useful in other low-resolution radiation-detection environments, including portable screening systems, drone-based radiation surveys, and waste-monitoring applications.

Several limitations should nevertheless be noted. The accuracy of the decomposition depends on the quality and stability of the reference spectra. Environmental changes, detector aging, gain drift, or altered background conditions could introduce bias if not properly accounted for. In addition, the present implementation assumes that the background used for compensation is representative of the measurement condition. Future work may therefore include adaptive background modeling, temporal background estimation, uncertainty-aware decomposition methods, and validation using a broader range of radionuclides and shielding scenarios.

In summary, the combination of BCCS preprocessing, complementary background compensation, and NNLS-based spectral decomposition provides an effective and interpretable framework for radionuclide identification using low-resolution plastic scintillation detectors. The method improves both visual separability and quantitative decomposition performance under weak-signal conditions and shows strong potential for practical deployment in next-generation RPM applications.


\section*{Methods}
\subsection*{Experimental Setup}

Gamma-ray spectra were measured using a large-volume plastic scintillation detector with dimensions of 30 cm$\times$180 cm$\times$5 cm, coupled to a photomultiplier tube and a multichannel analyzer with 2048 channels. Three radionuclides—Am-241, Cs-137, and Co-60—were used as reference sources. Each source spectrum was measured together with an independently acquired background spectrum at multiple source-to-detector distances.

To evaluate mixed-source conditions, nine combinations of the three radionuclides were also measured at different distances. All measurements were conducted under representative radiation portal monitoring conditions, and the acquisition times were selected to reflect practical low-count environments. As expected for plastic scintillation detectors, the raw spectra were


\subsection*{Spectral Channel Configuration}

In scintillation detector systems, spectral data are commonly acquired using 256, 512, or 1024 channels, whereas high-resolution semiconductor detectors such as HPGe often use much finer channel configurations, such as 4k or 8k channels. However, for plastic scintillation detectors, increasing the number of channels does not necessarily improve radionuclide discrimination because of the intrinsically low energy resolution of the detector response.

In this study, the spectral data were analyzed using 256 channels. The reason for this choice is that increasing the number of channels reduces the counts per channel and degrades counting statistics, while providing little additional discriminative information for plastic scintillator spectra. Therefore, the use of 256 channels provides a practical balance between spectral representation and statistical reliability for plastic scintillator-based measurements. This choice is particularly appropriate for plastic scintillation detectors, in which statistical
fluctuations are more influential than fine spectral resolution.


\subsection*{Spectral Preprocessing: Backward Cumulative Channel Sum (BCCS)}

The backward cumulative channel sum (BCCS) is defined as a cumulative transformation of the pulse-height spectrum. Let $\mathbf{C} = [C_1,\,C_2,\ldots,\,C_N]$ denote the raw pulse-height spectrum, where $C_i$ is the count in channel $i$. The BCCS spectrum is given by

\begin{equation}
B_i = \sum_{j=i}^N C_j. \label{eq:BCCS}
\end{equation}

This transformation produces a monotonically decreasing representation of the spectrum, reducing statistical fluctuations while preserving the global spectral shape. Because plastic scintillator spectra are dominated by broad Compton-related distributions rather than sharp photopeaks, preservation of the overall spectral shape is more important than retaining fine channel-by-channel variations.


During method development, both forward and backward cumulative summation schemes were examined. The forward cumulative sum is defined as

$$
F_i = \sum_{j=1}^i C_j. \label{eq_FCCS}
$$

In the forward cumulative representation, the cumulative values are constructed from the lowest channels upward. Because the spectral features relevant for radionuclide discrimination in plastic scintillators are primarily located in the low- to medium-channel region, the forward scheme provides limited noise suppression in these critical channels. As a result, statistical fluctuations in the most informative part of the spectrum remain largely unmitigated.
In contrast, the backward cumulative representation integrates counts from each channel toward the high-channel end of the spectrum. This leads to substantial averaging even in the low- and medium-channel region, where most of the physically meaningful spectral differences reside. Consequently, the backward scheme more effectively reduces statistical fluctuations in the most relevant part of the spectrum while preserving the global spectral trend.
For these reasons, the backward cumulative channel sum (BCCS) was adopted in this study as the primary spectral transformation method.


For each measurement condition, repeated acquisitions were performed more than 10 times under identical experimental settings. Because the run-to-run variation in spectral shape and in the resulting BCCS trend was negligible, a single representative spectrum is presented in the figures for clarity.

% Unlike the conventional forward cumulative spectrum that sums from the first channel upward, the BCCS accumulates counts from channel 
% $i$ to the end of the spectrum. This produces a smooth, monotonically decreasing cumulative profile that suppresses statistical fluctuations in the high-channel region while preserving the global spectral structure. Because this transform enhances separability under low-count conditions, it serves as the primary preprocessing step for all subsequent analyses.


\subsection*{Background Compensation Techniques}

After BCCS preprocessing, background compensation was applied to reduce the influence of ambient radiation and system background on the transformed spectra. In practical RPM environments, both the intensity and spectral shape of the background may vary depending on container structure, cargo composition, and measurement conditions. Therefore, appropriate background compensation is essential prior to radionuclide spectral decomposition.

Two background compensation techniques were considered in this study. The first was a subtraction-based method, defined as 

\begin{equation}
S^M_i =  B^M_i - B_i^{\textrm{bkg}}. \label{eq:DS}
\end{equation}

where $B_i^M$ and $B_i^\text{bkg}$ denote BCCS spectra of the measured signal and background, respectively. This approach preserves the net count scale of the source contribution and is suitable when absolute signal magnitude is relevant for interpretation, although it may also amplify statistical fluctuations when the count level is low.

The second was a normalization-based method, defined as 

\begin{equation}
R^M_i = \left(\dfrac{B^M_i}{B^{\textrm{bkg}}_i}\right) -1.\label{eq:NB}
\end{equation}

This formulation represents the relative deviation of the measured spectrum from the background. The subtraction of unity ensures that 
$R_i^M = 0$ corresponds to pure background, providing a convenient baseline for comparison. Because the ratio removes the common scaling factor between the measured and background spectra, it is largely independent of acquisition time and overall count scale. This property is particularly useful in real-time inspection systems, where the effective measurement time may vary as a container truck passes the detector.


In both compensation schemes, the spectral data can also be converted into count-rate form by dividing the measured counts by the acquisition time. In this case, the spectral values are no longer restricted to integers and may take fractional values. However, this does not affect the BCCS procedure, because the cumulative summation can be applied equally to integer and non-integer values. Therefore, the proposed method remains applicable even when the measurement time varies. This property is particularly important in practical RPM environments, where container trucks do not always pass the detector at a constant speed.

Both compensation methods were applied in the BCCS domain, and their performances were compared in terms of spectral contrast, stability, and robustness in subsequent decomposition analysis.

\subsection*{Radionuclide Activity Estimation via Linear Spectral Unfolding}

Following BCCS preprocessing and background compensation, the processed spectrum was analyzed using a linear spectral unfolding framework. Under fixed detector geometry and comparable background conditions, the processed spectrum can be approximated as a linear combination of radionuclide-specific reference spectra.

For the subtraction-based representation, the processed spectrum is expressed as
\begin{equation}
S_i^{M} = \sum_{X} \alpha_X S_i^{X},
\label{eqn:subtraction_unfolding}
\end{equation}
where $\alpha_X$ denotes the relative contribution of radionuclide $X$.

For the normalization-based representation, the situation is slightly different because the normalized spectrum depends explicitly on the background spectrum used for normalization. If the background spectra of the measured and reference data are sufficiently similar under the same measurement conditions, i.e.,
\begin{equation}
B_i^{M,\mathrm{bkg}} \approx B_i^{X,\mathrm{bkg}},
\label{eqn:bkg_similarity}
\end{equation}
the normalized processed spectrum can be approximated as
\begin{equation}
R_i^{M} \approx \sum_{X} \alpha_X R_i^{X}.
\label{eqn:ratio_unfolding}
\end{equation}

In both cases, the radionuclide contributions were estimated by solving a non-negative least squares (NNLS) problem:
\begin{equation}
\{\alpha_X\}
=
\argmin_{\alpha_X \ge 0}
\sum_{i=1}^{N}
\left(
I_i^{M} - \sum_{X}\alpha_X I_i^{X}
\right)^2,
\label{eqn:nnls}
\end{equation}

\noindent where $I_i^{M}$ denotes the processed spectrum of the measured data, $I_i^{X}$ denotes the processed reference spectrum for radionuclide $X$, and $I_i$ represents either the subtraction-based or normalization-based spectral form depending on the compensation method used. The non-negativity constraint ensures physically meaningful solutions by preventing negative radionuclide contributions.

Because meaningful source-dependent information in plastic scintillator spectra is concentrated mainly in the low- and medium-channel region, the unfolding analysis was performed over a restricted channel range. This reduces the influence of high-channel statistical fluctuations while preserving the spectral information relevant for radionuclide discrimination.

The resulting coefficients do not represent absolute radionuclide activities in becquerels, because the moving measurement geometry does not permit direct activity quantification. Instead, they are interpreted as reference-based indicators of relative contribution and relative signal level for each radionuclide.






\section*{Conclusion}

In this study, a BCCS-based spectral processing framework was proposed for radionuclide discrimination using large plastic scintillation detectors under practical RPM-like conditions. The method was designed to preserve the global spectral characteristics of plastic scintillator responses while reducing local statistical fluctuations through cumulative transformation.

The results showed that the BCCS representation provides a stable basis for comparing radionuclide-dependent spectral patterns and for decomposing mixed-source spectra in combination with background compensation and NNLS-based unfolding. The subtraction-based representation preserves net signal-scale information, whereas the normalization-based representation provides improved flexibility under varying acquisition conditions.

The unfolding coefficients obtained in this framework should be interpreted as reference-based indicators of relative contribution and relative signal level, rather than as absolute radionuclide activities. Within this interpretation, the proposed method provides a practical and physically meaningful approach for extracting radionuclide-specific information from low-resolution plastic scintillator spectra.

Although the present study focused on three representative radionuclides, the framework can be extended to additional radionuclides when appropriate reference spectra are available. The proposed approach is therefore expected to be useful for real-time radiation monitoring and container inspection applications.


\section*{Acknowledgements}

This research was supported by the Ministry of Oceans and Fisheries, Korea, through the Korea Institute of Marine Science \& Technology Promotion (KIMST) under Project No. 20200611. The authors acknowledge the use of ChatGPT (OpenAI) for language refinement and editorial assistance during the preparation of this manuscript. All scientific content, interpretations, and conclusions presented herein were solely developed and validated by the authors.

\section*{Author contributions statement}
M. Moon conceived the study and performed the experiments. J. So developed the analysis software and performed signal processing. B. Jeon synthesized the test datasets and supported evaluation. All authors contributed to writing and reviewing the manuscript.


\bibliography{biblio}


\begin{figure}[ht]
\centering
% \includegraphics[width=\linewidth]{figure/fig01_bkgref.pdf}
\includegraphics[width=\linewidth]{newfigures/figure1.pdf}
\caption{Spectral comparison of the background and reference radionuclides before and after BCCS processing. (a) Raw spectra of the background, Am-241, Cs-137, and Co-60 measured with the plastic scintillation detector. (b) Corresponding BCCS-transformed spectra. (c) Background-subtracted BCCS spectra. (d) Background-normalized BCCS spectra. The BCCS transformation produces smooth, monotonically decreasing profiles while preserving radionuclide-dependent spectral differences. Small Am-241 values in the high-channel region originate from statistical fluctuations in the measured low-count data and do not represent physically meaningful spectral structure.}

\label{fig:bkgref}
\end{figure}

\newpage 

\begin{figure}[ht]
\centering
% \includegraphics[width=\linewidth]{figure/fig04_CSBS.pdf}
\includegraphics[width=\linewidth]{newfigures/figure2.pdf}
\caption{Comparison of BCCS spectra for nine mixed-source configurations using two background compensation methods. (a) Background-subtracted BCCS spectra. (b) Background-normalized BCCS spectra. The subtraction-based representation preserves absolute net signal-scale differences among the mixed-source cases, whereas the normalization-based representation expresses the spectra as relative deviations from the background and provides improved comparability under varying acquisition conditions.}
\label{fig:mixedbccs}
\end{figure}

\newpage


\begin{figure}[ht]
\centering
\includegraphics[width=\linewidth]{newfigures/fit_BGSub.pdf}
\caption{Linear spectral unfolding results for mixed-source spectra using subtraction-based BCCS processing. The processed BCCS spectra (solid lines) are fitted by linear combinations of the reference radionuclide spectra (dashed lines). The results show that the mixed-source spectra can be decomposed into Am-241, Cs-137, and Co-60 components with non-negative coefficients. Small residual contributions may also appear for radionuclides not intentionally included in a given source configuration because of limited separability in low-resolution plastic scintillator spectra and statistical fluctuations.}
\label{fig:fi_fitbgsub}
\end{figure}

\begin{figure}[ht]
\centering
\includegraphics[width=\linewidth]{newfigures/fit_BGRto.pdf}
\caption{Linear spectral unfolding results for mixed-source spectra using normalization-based BCCS processing. The normalized processed spectra are fitted by linear combinations of the normalized reference spectra. The normalization-based representation provides stable decomposition while reducing sensitivity to variations in overall count scale and effective acquisition time. As in the subtraction-based case, small residual contributions may appear for radionuclides not present in a given source configuration.}
\label{fig:fi_fitbgrto}
\end{figure}

\begin{figure}[ht]
\centering
\includegraphics[width=\linewidth]{newfigures/fit_contributions.pdf}
\caption{Variation of the estimated radionuclide coefficients with source-to-detector distance for the subtraction-based and normalization-based methods. The coefficients obtained from spectral unfolding represent reference-based relative signal levels rather than absolute activities. The horizontal dashed lines indicate, for each radionuclide, the mean coefficient obtained from NNLS fitting of spectra in which that radionuclide was absent, and serve as practical reference levels for discrimination. As the distance increases, the estimated coefficient decreases because of geometric attenuation and the corresponding reduction in net detector response. The plotted results are representative; each condition was repeated more than 10 times under identical settings, and the overall trend was reproducible.}
\label{fig:fi_fitcontrib}
\end{figure}

\end{document}